\subsection{Spectrum of a ring}
\begin{exer}
	For $k$ a field, $\Spec (k[[T]])$ is the space $\{ \xi, 0 \}$ with
	generic point $\xi$ and closed point $0$.
\end{exer}
\begin{proof}
	It is easy to see that any power series with a non-zero constant term is
	a unit. Furthermore, if we have $f = \sum a_nT^n$ with $a_N = 1$ for
	some $N$ and $a_n = 0$ whenever $n \leq N$, then we can pull out a
	common factor of $T^N$. We see that $f$ is then a unit multiple of $T^N$
	and $(f) = (T^N)$. Clearly, this is only prime when $N = 1$. It is easy
	to see that $k[[T]]$ is an entire ring, so the proposition is immediate.
\end{proof}

\begin{exer}
	Let $\varphi: A \to B$ be a homomorphism of finitely generated algebras
	over a field $k$. The image of a closed point under $\Spec \varphi$ is 
	a closed point.
\end{exer}
\begin{proof}
	If $\Spec \varphi$ takes the maximal ideal
	$\mathfrak m$ into the ideal $\mathfrak p$, we get a diagram
	\[
		k \xrightarrow{} A/\mathfrak p \xrightarrow{\widetilde{\varphi}}
		B/\mathfrak m.
	\]
	By Corollary 1.12, $B/\mathfrak m$ is a finite algebraic extension of
	$k$, so $B/\mathfrak m$ (resp. $A/\mathfrak p$) is finite over
	$A/\mathfrak p$ (resp. $k$). It then suffices to prove that
	$A/\mathfrak p$ is a field. Any $a\in A$ has an associated
	irreducible polynomial $P = \sum a_i X^i \in k[x]$ such that $P(a) = 0$.
	This polynomial must have a nonzero constant term, so we have
	\[
		-\frac {1}{a_0} \sum_{i > 0} a_i a^i = 1,
	\]
	but the polynomial on the left has no constant term and we can pull out
	a factor of $a$.
\end{proof}

\begin{exer}
	Let $k = \RR$ be the field of real numbers and let
	$A = k[X,Y]/(X^2 + Y^2 + 1)$. Let $x,y$ be the respective images of
	$X,Y$ in $A$.
	\begin{enumerate}[label=(\alph*)]
		\item Let $\mathfrak m$ be a maximal ideal of $A$. Then there
			are $a,b,c,d \in k$ such that
			$x^2 + ax + b,y^2 + cy + d \in \mathfrak m$.
			Furhtermore, $\mathfrak m$ is generated by a single
			element $f = \alpha x + \beta y + \gamma$ with at least
			one of $\alpha,\beta$ nonzero.
		\item The map $(\alpha,\beta,\gamma) 
			\mapsto (\alpha x + \beta y + \gamma)A$ establishes a
			bijection between the subset
			$\PP(k^3) \setminus \{ (0,0,1) \}$ of the projective
			space $\PP(k^3)$ and the set of maximal ideals of $A$.
		\item Let $\mathfrak p$ be a non-maximal prime ideal of $A$.
			Then the canonical homomorphism $\varphi: k[X] \to A$
			is finite and injective, and
			$k[X] \cap \mathfrak p = 0.$ Let $g \in \mathfrak p$
			and let $g^n + a_{n-1}g^{n-1} + \ldots + a_0 = 0$ be an
			integral equation with $a_i \in k[X]$. Then
			$a_0 = 0$, and so $\mathfrak p = 0$.
	\end{enumerate}
\end{exer}
\begin{proof}
	First, notice that $A/\mathfrak m$ is isomorphic to $\CC$, since
	$X^2 + Y^2 + 1$ does not split in $k$. Then the existence of $a,b,c,d$
	follows from the fact that any complex number satisfies a monic
	quadratic equation. Since $\CC$ is degree 2 over $k$, there are
	$\alpha,\beta,\gamma$, not all zero, such that
	\[
		\alpha x' + \beta y' + \gamma = 0,
	\]
	where $x'$ and $y'$ are the respective images of $x$ and $y$. Since not
	all of them are zero, we may assume that $\alpha$ is nonzero. Then we
	get an equation for $x'$ in terms of $y'$, which we can then substitute
	into the equation $x'^2 + y'^2 + 1 = 0$. This new equation cannot split
	in $k$ because that would imply a real root of $x'^2 + y'^2 + 1 = 0$,
	so we have that $\mathfrak m = fA$, where $f$ is the ideal
	corresponding to the linear equation in $x'$ and $y'$.\\

	Part (b) holds, since a solution to $X^2 + Y^2 + 1 =0$ is
	determined (up to conjugacy) by an arbitrary complex number, which
	itself determines a linear equation with real coefficients.\\

	That $\varphi$ is finite and injective is clear, so consider the
	induced homomorphism
	$\widetilde{\varphi}: k[X]/\varphi^{-1}(\mathfrak p)\to A/\mathfrak p.$
	Notice that $B = k[X]/\varphi^{-1}(\mathfrak p)$ is not a field, since
	$\widetilde{\varphi}$ is still finite and injective, which implies that
	$A/\mathfrak p$ is a field as well. Then $B = k[X]$ and
	$\mathfrak p \cap k[X] = 0$. Finally, in the integral equation for $g$,
	notice that $a_0$ must be a multiple of $g$, hence 0. Then the equation
	has a common factor of $g$ which we may cancel, which gives a new 
	equation with constant term $a_1$. This repeats indefinitely and so
	$g = 0$.
\end{proof}

\begin{exer}
	Let $A$ be a ring.
	\begin{enumerate}[label=(\alph*)]
		\item Let $\mathfrak p$ be a minimal prime ideal of $A$. Then
			$\mathfrak pA_{\mathfrak p}$ is nilpotent. In
			particular, every element of $\mathfrak p$ is a
			zero divisor.
		\item If $A$ is reduced, then any zero divisor in $A$ is an
			element of a minimal prime ideal. This is not true in
			general if $A$ is not reduced.
	\end{enumerate}
\end{exer}
\begin{proof}
	For (a), notice that $I = \mathfrak pA_\mathfrak p$ is the only prime
	ideal in $B = A_\mathfrak p$, so $I = \Nil(B)$. It is easy to see that
	the radical of any principal ideal $(x) \subseteq B$ with $x \in I$ is
	itself just $I$, so the powers of any element $x \in I$ give $I$ and
	$I$ is finite. The second part of (a) is obvious.\\

	For (b), let $x$ be any zero divisor and let $\mathfrak p$ be a prime
	ideal minimal among those that contain $x$. If $\mathfrak p$ is minimal,
	we are done. Otherwise, suppose that $\mathfrak q \subset \mathfrak p$
	is a minimal prime ideal not containing $x$. Since $A$ is reduced, we
	must have $\sqrt{(x)} \cap \Ann(x) = \{ 0 \}$, so $\mathfrak q$ must
	contain a nonzero element $x' \in \Ann(x)$. Localizing by $\mathfrak q$,
	we see that $x$ is sent to $0$ and $x'$ is sent to a nilpotent. This
	implies that there is some positive integer $n$ and some element
	$s \in A \setminus \mathfrak q$ such that $s(x'^n-x) = 0$, and it is
	easy to see that this means that $sx^2 = 0 \in \mathfrak q$, so 
	$x^2 \in \mathfrak q$, contradiction.
\end{proof}
The fact that this is false for $A$ not reduced comes is immediate from the fact
that there can be non-nilpotent zero divisors in a ring.\\

